\intro
За последние годы 3D-печать перешла от задач прототипирования к практическому применению в малосерийном производстве, учебных лабораториях и инженерных мастерских. Это стало возможным благодаря экономичности аддитивного подхода: он не требует оснастки при небольших партиях, уменьшает объём отходов по сравнению с традиционными методами и сокращает время изготовления изделий. Однако при переходе от единичных прототипов к серийному производству ключевым ограничением становится стабильность качества и точность результата.

В рамках данной работы рассматривается метод послойного наплавления (Fused Deposition Modeling, FDM). Изделие формируется последовательным нанесением расплавленного материала по слоям в соответствии с управляющей программой. Для FDM-процесса характерна высокая чувствительность к параметрам температуры, скорости печати и перемещений, охлаждения, адгезии первого слоя и к состоянию расходного материала. Отклонение даже одного из параметров может вызвать каскадное ухудшение качества изделия. Поэтому контроль процесса должен выполняться не постфактум, а непосредственно во время печати, когда еще возможно своевременно повлиять на ход задания.

Значимая часть отказов проявляется уже на ранних этапах печати и связана с нарушением устойчивости технологического процесса, что влияет на качество формируемого изделия. Если такие отклонения не обнаруживаются своевременно, последствиями становятся перерасход материала, потеря машинного времени, необходимость повторного запуска задания и рост эксплуатационной нагрузки на узлы принтера.

Дополнительной проблемой является простой оборудования между заданиями. На фермах 3D-принтеров и в учебно-производственных лабораториях операторы часто работают по графику~5/2, что формирует ночные и выходные периоды неполной загрузки. Частичный переход к дневным и ночным сменам позволяет сократить простой, но одновременно увеличивает затраты на персонал. Следовательно, для повышения общей эффективности необходимо снижать зависимость производственного цикла от постоянного присутствия оператора и переводить контроль печати в автоматизированный режим.

В рамках исследования разрабатывается 3D-принтер с функцией экспериментальной платформы и контуром раннего обнаружения дефектов во время печати. Этот контур включает видеонаблюдение зоны печати, обработку кадров, выявление и локализацию дефектов по изображению, а также передачу данных в управляющую логику принтера. Под реагированием подразумеваются заранее заданные действия, направленные на ограничение развития дефекта и снижение потерь материала и времени.

Полная автоматизация цикла, включая расширенную очередь заданий и автоматическое снятие изделий, рассматривается как перспектива дальнейшего развития. Однако в данной работе это не является основной целью. Тем не менее, основы для этого будут заложены уже сейчас.

Объектом исследования является процесс FDM-печати в условиях малосерийного производства и длительной эксплуатации оборудования. Предметом исследования являются методы и программно-аппаратные средства раннего обнаружения дефектов 3D-печати по данным видеоконтроля, а также способы их интеграции в контур управления принтером для оперативного принятия решений.

Практическая значимость работы состоит в том, что полученная первая версия системы задаёт основу для перехода от ручного контроля к контролю в процессе печати. Ранняя фиксация отклонений позволяет сократить потери материала и времени, а также формирует архитектуру, пригодную для дальнейшего расширения функций автономной работы.

Актуальность работы определяется тем, что небольшие нарушения режима FDM-печати приводят к браку и повторным запускам, увеличивая потери материала и времени и снижая загрузку оборудования. В таких условиях автоматическое обнаружение и локализация дефектов с реакцией в ходе печати является практической и научно-технической задачей. На момент написания работы в открытых источниках и на российском рынке не выявлено распространённого решения для FDM, которое встраивается в контур управления принтером и обеспечивает раннее обнаружение дефектов по видеопотоку с заранее заданными сценариями реакции. При этом существуют зарубежные системы мониторинга со схожей целью (в том числе с функциями обнаружения отказов с применением искусственного интеллекта (ИИ) и действиями вроде уведомления или паузы печати), однако их принципы работы и степень интеграции с управляющей логикой отличаются от рассматриваемого в данной работе подхода. Поэтому разработка и экспериментальная проверка первой версии системы для FDM-печати с функцией раннего обнаружения дефектов и реакцией на них является актуальной задачей, которая может внести вклад в повышение эффективности и качества 3D-печати.

Целью данной работы является разработка и экспериментальная проверка первой версии 3D-принтера с системой раннего обнаружения дефектов в процессе печати и механизмом своевременного реагирования. Для достижения поставленной цели необходимо решить следующие задачи:
\begin{enumerate}
\item проанализировать технологический цикл FDM-печати и выделить дефекты, критичные для раннего обнаружения;
\item сформировать требования к архитектуре программно-аппаратного контура, включающего принтер, камеру, обработку данных и подсистему управления;
\item организовать подготовку и разметку визуальных данных для обучения и проверки модели обнаружения дефектов;
\item разработать модель обнаружения дефектов для целевых классов;
\item реализовать интеграцию модуля обнаружения в рабочий цикл принтера и настроить сценарии своевременного реагирования на выявленные отклонения;
\item провести экспериментальную проверку работоспособности первой версии на серии печатей и зафиксировать ограничения текущей реализации;
\item определить направления развития решения, включая расширение автономности цикла и развитие очереди заданий.
\end{enumerate}
